\hypertarget{appendix-91d09ebed979}{}
\section{Serialization}\label{appendix-91d09ebed979-serialization}

\textbf{Curriculum depth:} Developed · \textbf{First seen:} C001 · \textbf{Latest development:} C006

\subsection{Definition and mechanism}\label{appendix-91d09ebed979-definition-and-mechanism}

Serialization maps semantic fields to a specified byte sequence; decoding validates and reconstructs them. Width, order, version, bounds, and padding policy belong to that contract. A U8/U32-LE/U64-LE header occupies 13 bytes at offsets 0/1/5, independently of a native record's possible 16-byte ABI layout. Packing removes some holes, not host dependencies. Publication and durability are subsequent operations.

M004 extension: a binary32 representation word such as 3F800000 can be serialized as little-endian bytes 00 00 80 3F. Raw-bit preservation retains signed-zero and NaN-payload differences; canonicalization deliberately merges them. Numeric equality does not imply byte identity. This is a format policy, not permission to change existing files.

M006 adds explicit schema/version/identity contracts, overflow-safe dimensions and resource limits, raw-word payload preservation, integrity coverage and a complete 52-byte demonstration. A native scalar writer is not automatically cross-endian stable.

\subsection{Prerequisites}\label{appendix-91d09ebed979-prerequisites}

\begin{itemize}
\tightlist
\item
  \hyperref[appendix-f414bd4ed05c]{Interpretation contract} --- Representation, type, and lifetime must identify what object exists.
\item
  \hyperref[appendix-cc04044b5782]{Fixed-width integer types} --- The representation width fixes byte count and valid shift range.
\item
  \hyperref[appendix-bc07c3d25b18]{Endianness} --- Byte order tells how digit significance maps to sequence positions.
\end{itemize}

\subsection{Used by}\label{appendix-91d09ebed979-used-by}

\begin{itemize}
\item
  \hyperref[appendix-590bbb2424c1]{Binary format contracts} --- uses this mechanism to define its representation, error or validation contract.
\item
  \hyperref[appendix-ed338677de0a]{Checksums CRCs and hashes} --- uses this mechanism to define its representation, error or validation contract.
\item
  \hyperref[appendix-e90053a4b201]{Floating-point canonicalization} --- uses serialization to establish this mechanism's representation or policy.
\item
  \hyperref[appendix-e89a8f290ef5]{Atomic publication and durability} --- uses this mechanism as stated in its prerequisites.
\end{itemize}

\subsection{Recall and next depth}\label{appendix-91d09ebed979-recall-and-next-depth}

Explain why packed/native memory dumps do not establish a portable schema.

\textbf{Pending:} Version migration, malformed-input protocols, other floating formats and cross-platform special-value transport.

\subsection{Source and navigation}\label{appendix-91d09ebed979-source-and-navigation}

\hyperref[01-foundations-i-ch002-bytes-become-a-durable-agreement]{Narrative development}

\hyperref[01-foundations-i-ch001-a-stretching-ruler-for-real-valued-quantities]{C004 representation policy}

\hyperref[01-foundations-i-ch002-where-the-bytes-live]{C003 teaching chapter} · \hyperref[appendix-fd090d982a15-m003---parts-41-60]{Part definitions} · \href{../02-Foundations-II/main.pdf\#nameddest=appendix-4ad47bdd269a}{Dependency Map (02-Foundations-II, appendix-4ad47bdd269a)}
